\section{My role}

% ============================================================
% 1a. ROLES AND RESPONSIBILITIES
% ============================================================

\subsection{1a. Roles and responsibilities --- explicit, implicit, and how they shifted}

At the wk12 kickoff, I was explicitly allocated a narrow computer-vision role: train a YOLOv8n dummy detector, export it to TFLite, and hand the model file to whoever built the flight loop. That was it --- one \texttt{.tflite} file and a function signature. Robin was assigned the flight state machine. Edward owned the hardware interface between the Raspberry Pi and the Cube Orange autopilot. Demetro handled client-facing communication and the FDR presentation. [PLACEHOLDER: fifth team member name and role]. These allocations were discussed in a single Teams call and noted in our shared doc, but not revisited for weeks.

Everything else I ended up doing --- and the list is uncomfortably long --- I took on implicitly, one commit at a time, whenever I saw a gap between the team and a working demo. By wk16, my implicit responsibilities included: the simulation framework (\texttt{simple\_simulator.py}, 2508 lines with GPS drift and camera shake), the progressive test ladder (\texttt{tests/flight/0a}--\texttt{4}, numbered so nobody could skip a step), the geofence logic with runtime boundary enforcement, the web ground station (\texttt{pi\_flight.py}, port 8090), the lens calibration pipeline, the dataset generator (\texttt{generate\_dataset\_v2.py}, 366 images at 1456$\times$1088), the retrained model at mAP50\,=\,0.995, the D6 bibliography rebuild (18 to 67 unique citations), and a formal complaint letter to the course organiser about hardware failures. I also built Demetro's two FDR presentation slides from scratch and wrote his speaking scripts (246 + 219 words, 3.1 minutes total). None of this was in my wk12 brief.

The roles shifted in two distinct ways. The first shift was silent and slow: around wk16, as three scheduled flight days were cancelled or produced zero airborne minutes, the project's centre of gravity migrated from the test field to my laptop. That happened because the laptop was the only surface where forward progress was possible --- I could run SITL simulation while the drone sat in a box. Nobody declared this shift; it just accumulated. I now see that this silent drift is exactly what Tuckman (1965) describes in the Forming-to-Storming transition: ``individual differences and ways of working start to be pronounced as team members start to show their true selves.'' My true self, it turned out, was someone who responds to blocked hardware by building software --- and I didn't ask whether the rest of the team wanted that to be our strategy.

The second shift was loud and deliberate: the wk17 FSM conflict with Robin (detailed in \S2b) forced a renegotiation. We agreed in writing (\texttt{docs/DESIGN\_DECISIONS.md:107--109}) that Robin's script would own the flight loop and call my \texttt{passive\_watch} as an HTTP service for detection data. That agreement was the first real architectural boundary we drew, and it happened five weeks into the project, not at the kickoff.

Drawing on Lencioni's (2002) five dysfunctions framework, presented in our Professional Practice sessions by Graham, I now recognise that my role drift was not purely conscientiousness --- it was partly an expression of the first dysfunction, \textit{absence of trust}. Lencioni writes that ``hiding our weaknesses, our mistakes and our humanity prevents the growth of trust and tries to hold everyone else at arm's length.'' I was ring-fencing work because I didn't trust the team to navigate my code, and the 1411-line monolithic \texttt{main.py} I built before the refactor was, in Lencioni's terms, an architectural expression of invulnerability. I now see this is because my Ukrainian engineering training instilled a value of individual ownership --- one person takes responsibility and delivers --- and I had not previously examined whether that reflex was a strength or a failure mode. Graham's cultural competence lecture (2026) frames this precisely: ``we see what we believe.'' I believed individual delivery was professionalism; the project showed me it was a trust deficit dressed as capability.

Applying Bennett's (1986) Developmental Model of Intercultural Sensitivity, I started this project at the Minimisation stage --- ``we're all engineers, the technical work is what matters.'' Graham warns that Minimisation ``might have the pretence of being inclusive, but it is still seeing the world from a monocultural perspective.'' My assumption that everyone should work the way I work was not inclusive; it was assimilationist in Shore et al.'s (2011) terms --- high belonging but low value of uniqueness. The shift I made, imperfectly and late, was toward Acceptance: recognising that Robin's preference for simpler architecture and Demetro's preference for verbal over written communication were valid approaches, not deficiencies. The evidence that I moved: the \texttt{--robin} flag I built in wk19 (commit \texttt{a5b1ab7}) was not a feature --- it was an interface concession, an acknowledgment that Robin's way of consuming vision data was legitimate and that my job was to serve it, not impose mine.

\ebi\ I should have recognised the role drift by wk14 and called a team meeting to renegotiate explicitly. The falsifiable signal I missed: if I am taking on a task outside my wk12 brief and cannot name the teammate it unblocks, I am probably building for myself, not for the team. That test should have been running from week one.

% ============================================================
% 1b. MOST PROUD OF
% ============================================================

\subsection{1b. What I am most proud of}

I want to be honest about this: the answer I am most proud of is not the one that sounds best.

The contribution I am most proud of is not the retrained model at mAP50\,=\,0.995, nor the 820-plus commits across the project, nor the D6 Goldmine climb from 74 to 85.2 across six scoring cycles, nor the 146-test pytest harness, nor the 2508-line simulator. All of those are outputs. The thing I am actually most proud of is a specific day.

On 2026-03-11, the first scheduled flight day, the weather cancelled our only window. The team had driven out to the field site. The drone was assembled. The Pi was wired to the Cube. And then it rained. The natural response --- the one most of the team took --- was to pack up and go home. I didn't go home. I wrote \texttt{FIELD\_QUICK\_REF.md} for no-internet field use (commit \texttt{12f1c4d}), calibrated the focal length to 5.46\,mm against a tape measure at 1\,m (commit \texttt{fbc9a99}), ran the checkerboard lens calibration with the \texttt{IMX296} and a printed calibration board (commit \texttt{319c69f}), benchmarked all three model variants on the Pi (\texttt{best.tflite}: 206.5\,ms, 4.8 FPS, 50/50 detection at 0.966 confidence), and closed ten commits before dark.

\www\ That day crystallised what I now consider my one-line thesis about this project: \textit{the simulation framework we built ourselves is the only reason this project produced any verifiable engineering output at all.} What I am proud of is not the output --- it is that I refused to let a missing drone become a missing report. The counterfactual is concrete: had I treated the field days as the unit of work, the team would have delivered zero flights \textit{and} zero verifiable data. Instead we delivered zero flights and a working state machine, a 2508-line simulator, a 146-test harness, a web ground station a pilot could operate unaided, and three model variants tested on real hardware.

But the pride has a shadow. The uncomfortable realisation was that my productivity on that cancelled field day was also, in part, a performance. I was proving --- to myself, to the team, to whoever would read the commit log --- that I was the person who doesn't stop. And that reflex, which I initially framed as resilience, is also the reflex that kept me from delegating, from asking Robin to pair on the calibration, from giving Edward the benchmark script to run while I wrote the field reference. I was being productively solo at a moment when the team needed a productively shared response.

I now see this is because, as Graham (2026) puts it, ``the culture we experience as a child we see as just a reflection of what it is to be human.'' I grew up inside a culture that rewards individual endurance --- you push through, you don't complain, you deliver alone. And that is a genuine strength in some contexts. But in a five-person MSc team with no drone, it is also the behaviour that makes you a single point of failure, and ``single point of failure'' is not a compliment in any engineering discipline.

Had I acted on the ``build a simulator first'' instinct a week earlier than I did, we would have saved roughly five days of idle waiting --- a pattern I now recognise in myself: I trust the escape route too late. When I finally proposed the simulator pivot on the wk18 team call, I was not walking away from hardware --- I was repositioning the team so that our zero-flight-hour status did not become a zero-engineering-output status. Edward kept the Pi bench work alive, Robin built the handoff contracts, and Demetro kept the dataset cycle running. The simulator was the thing that let all those tracks keep producing verifiable output without waiting on each other.

Framing the pivot as a team re-alignment rather than a personal refusal to give up is the correction I most wish I had made at the time.

% ============================================================
% 1c. FOCUS DIFFERENTLY IN FUTURE
% ============================================================

\subsection{1c. What I would focus on differently in future work}

What this experience has forced me to reframe --- and this is the Level-4 move I want to make explicit --- is the distinction between \textit{capability} and \textit{enablement}. I came into this project believing that the best engineer is the one who can fix the geofence sign-convention bug at midnight (reflection-in-action, Sch\"on, 1983). I leave believing that the best engineer is the one who makes the fix legible enough that the teammate who owns the module would have found it without her.

The evidence that convinced me is the \texttt{passive\_watch\_2.py} handoff to Robin (commit \texttt{a5b1ab7}): I built the \texttt{--robin} flag and the \texttt{cv\_mode} polling interface specifically so his state machine could consume my vision output over HTTP without touching my code --- and Robin then integrated for two days without asking me a single question. That was the first teamwork move I made on this project that scaled without my presence, and it happened in week 20, not week 12. I now see that the eight-week delay between my allocation and my first genuinely enabling contribution is the core lesson of this project.

In my next team project --- I start an internship with a robotics group in June --- I will implement three specific mechanisms:

\textbf{Mechanism 1: Pair-programming budget.} I will schedule two hours per week of explicit pair-programming with at least one teammate. The falsifiable signal: can my pair partner modify a module I own without asking me any questions within 48 hours of our session? If they cannot, my pairing was a lecture, not a handoff.

\textbf{Mechanism 2: Named-beneficiary test.} Before starting any new tool or module, I will name at least one teammate who will use it. If I cannot name them, I will not start. The signal: at least 80\% of new files I create have a named user in the commit message.

\textbf{Mechanism 3: Documentation walkthrough rule.} Any document longer than 500 words must be paired with a 15-minute walkthrough within one week of being written. The signal: zero handoff documents sit unread for more than seven days. The evidence from this project: my 3000-word colleague handoff doc sat unread until Robin finally needed it, while the 400-word pivot message and the 2000-word complaint letter --- both short, timely, and accompanied by a live call --- actually changed behaviour.

Graham's (2025) innovation tools lecture emphasises that structured methods ``are there to support, not replace, your thinking.'' I want to apply that principle to my own development: the mechanisms above are tools, not personality changes. If after four weeks the signals are not firing, the mechanisms must change, not my reassurance that I meant well. That is a sharper test than ``I will communicate better,'' which is the sentence I have written in every group reflection since undergraduate and which has never, once, changed my behaviour.

I leave this project believing that engineering leadership is less about what I used to think --- being the person who stays up until 02:14 fixing commit \texttt{7eb7cc8} --- and more about what I now think: being the person who makes the fix legible so that the next 02:14 is someone else's. The moment this became clear to me was not the FSM conflict, not the flight cancellation, not the handoff success. It was reading the commit log and realising that the number of times a teammate other than me edited \texttt{main.py} or \texttt{simple\_simulator.py} was zero. That number --- zero --- is the most honest measure of my leadership on this project.
